Disturbances need to have a continuity.


\section{\index{CONTINUE SOLVER}Continue Solver}\label{sec-continue-solver}

time(s) CONTINUE SOLVER disc\_meth max\_h(s)
	min\_h(s) latency(pu) upd\_over  
	
Discretization method (disc\_meth): 
\begin{itemize}
	\item TR: Trapezoidal 
	\item BE: Backward Euler 
	\item BD: BDF2 
\end{itemize}
Jacobian update override (upd\_over):
\begin{itemize}
	\item ALL: force a full Jacobian refactorisation
	\item NET: force a refactorisation of the network Jacobian
	\item ABL: force an update of the injector Jacobian at every bus
	\item IBL: force nothing in addition, but keep the bus Jacobian updates already requested by the other disturbances applied at this time instant
	\item NOT: override, discarding even those; the bus Jacobian flags are restored to their values on entry
\end{itemize}
It is mainly used to modify the settings of the solver and has to
exist at the first line. For example:

0.000 CONTINUE SOLVER BD 0.0200 0.001 0. ALL


\section{\index{CONTINUE DISPLAY}Continue Display}

time(s) CONTINUE DISPLAY new\_plot\_step(s)

Changes the sampling time step used for the output of observed
variables from that point of the simulation onwards. For example,
reducing the output sampling to one point per second after $t=20$~s:

20.000 CONTINUE DISPLAY 1.0


\section{Set stopping criteria for voltage}

Define in the data files the following record:

 DCTL SIM\_MINMAXVOLT CTRL\_Name VMAX(pu) VMIN(pu)
	DEADTIME(s) Stop\_Simulation(T/F) ;



\section{Set stopping criteria for machine speed}

Define in the data files the following record:

DCTL SIM\_MINMAXSPEED CTRL\_Name MAX\_SPEED(pu)
	MIN\_SPEED(pu) DEADTIME(s) Stop\_Simulation(T/F) ;


\section{\index{STOP}Stop}

time(s) STOP 

It is used to signal the end of the simulation and has to exist at the last line. 

For example:

100.000 STOP 

\section{\index{BREAKER}Trip line}

time(s) BREAKER BRANCH name\_of\_line orig\_break(0/1)
	extrem\_break(0/1) 
	
Used to open/close the breakers of a line. 

For example, opening both ends of a line:

10.000 BREAKER BRANCH 1044-4032 0 0 

\section{\index{BREAKER}Trip synchronous machine / injector}

time(s) BREAKER SYNC\_MACH/INJ name\_of\_synch\_mach/name\_of\_injector
	breaker(0/1) 
	
Used to open/close the breaker (trip) of a synchronous machine or injector. 

For example:

10.000 BREAKER INJ L\_11 0

\section{\index{BREAKER}Trip two-port}

time(s) BREAKER TWOP name\_of\_twoport 0

Used to disconnect a two-port component (e.g.\ an HVDC link) at both
ends. Only opening is supported, so the breaker status must be 0.

For example:

10.000 BREAKER TWOP hvdc1 0

\section{\index{FAULT}Three phase short-circuit (with use of resistance to
	ground)}

This demands two commands:

time(s) FAULT BUS name\_of\_bus rfault {[}xfault{]}
	
time(s) CLEAR BUS name\_of\_bus

The first line is used to declare the starting of the short-circuit.
The fault has a resistance of \emph{rfault+j{*}xfault} to the ground
where the values of rfault and xfault are in Ohm. If xfault is not
defined, a fully resistive fault is assumed.

Example of a 100ms short-circuit directly to ground:

10.000 FAULT BUS 1044 0. 0.\\
10.100 CLEAR BUS 1044
	

\section{\index{VFAULT}Three phase short-circuit (with use of voltage reached
	after fault)}

This demands two commands:

time(s) VFAULT BUS name\_of\_bus Voltage\_reached\_after\_fault
	
time(s) CLEAR BUS name\_of\_bus

The first line is used to declare the starting of the short-circuit.
Internally, RAMSES creates a temporary shunt admittance $B_{fault}$
at the faulted bus and iteratively adjusts it using a secant method
until the bus voltage matches the target value (declared in pu) to
within $\pm 0.5\,\%$. The injected currents are $i_x = B_{fault}\,v_x$
and $i_y = -B_{fault}\,v_y$. Up to 10 correction steps are attempted;
if convergence is not reached, the simulation halts with an error
message.

Example of a 100ms short-circuit where the voltage at the faulted
bus reached 0.5 pu after the fault:

10.000 VFAULT BUS 1044 0.5\\
10.100 CLEAR BUS 1044
	

\textbf{Supported:} Version 3.13 and above.


\section{\index{CHGPRM}Change parameters}\label{sec-chgprm}

Used to change the parameters of a model during the simulation.


\subsection{BRANCH}

time(s) CHGPRM BRANCH name\_of\_line MAGN/PHAN
	$\pm$increment 

\subsection{SHUNT}

time(s) CHGPRM SHUNT name\_of\_shunt QNOM $\pm$increment

The increment should be expressed in MVAr and it
is per unitized internally by RAMSES.


\subsection{EXC}

time(s) CHGPRM EXC name\_of\_equipment name\_of\_parameter
	$\pm$increment (MVAr/\%) duration(s) 
	
The units is
not obligatory. If nothing is given then the parameter is modified
in absolute value. If MVAr is given, then the increment is per unitized
using \emph{Snom} of the machine before applying the change. If \%
is given, then the parameter is changed as a percentage of the original
value. If $duration=0$ then a step change is applied, otherwise the
change is applied as a ramp over the given duration (in seconds).

For example:

10.000 CHGPRM EXC g1 V0 +10 \% 10 

This
means the parameter V0 of the exciter of synchronous machine g1 is
ramped by +10\% between 10 and 20 seconds.


\subsection{TOR}

time(s) CHGPRM TOR name\_of\_equipment name\_of\_parameter
	$\pm$increment (MW/\%) duration(s) 
	
The units is
not obligatory. If nothing is given then the parameter is modified
in absolute value. If MW is given, then the increment is per unitized
using \emph{Pnom} of the machine before applying the change. If \%
is given, then the parameter is changed as a percentage of the original
value. If $duration=0$ then a step change is applied, otherwise the
change is applied as a ramp over the given duration (in seconds). 

For example:

10.000 CHGPRM TOR g1 P0 +1 MW 10 

This
means the parameter P0 of the torque controller of synchronous machine
g1 is ramped by +1 MW between 10 and 20 seconds.


\subsection{INJ/TWOP/DCTL}

time(s) CHGPRM INJ/TWOP/DCTL name\_of\_equipment
	name\_of\_parameter $\pm$increment (MW/MVAr/\%/SETP) duration(s)

The units is not obligatory. If nothing is given
then the parameter is modified in absolute value. If MW or MVAr is
given, then the increment is per unitized using system's \emph{Sbase}
before applying the change. If \% is given, then the parameter is
changed as a percentage of the original value. IF SETP is given then
the value increment is actually the new setpoint. If $duration=0$
then a step change is applied, otherwise the change is applied as
a ramp over the given duration (in seconds). 

For example:

10.000 CHGPRM INJ L\_11 P0 +50 \% 60 \\
10.000 CHGPRM INJ L\_11 Q0 +30 \% 60 
	
	This means the
parameter P0 (resp. Q0) of the injector L\_11 is ramped by +50\% (resp.
30\%) between 10 and 70 seconds. In this case, if L\_11 is a load
model, it can be used to simulate a load increase.


\section{\index{JAC}Export Jacobian matrix}

time(s) JAC 'name\_of\_filename' 

Also, make sure to add these to your settings:

\$OMEGA\_REF SYN ;
	
\$SCHEME IN;


\section{\index{LFRESV}Export load flow}

Takes a snapshot of the system and exports the load flow at a specific time.

time(s) LFRESV 'name\_of\_filename' 